\section{Introduction}
\label{sec:intro}

\subsection{The problem}

A speedcubing result is a claim: \emph{this person took a cube in this
scrambled state and returned it to solved, unaided, in this many seconds}.
In person, a judge validates the claim by watching. Online, there is no
judge, and the natural substitute --- a smart cube that reports its own
turns over Bluetooth --- validates nothing, because the thing reporting the
turns is the thing under suspicion. A cube that can send a move stream can
send a fabricated one.

So the claim has to be decided from something the competitor does not
control. The only such thing available to a consumer is the video of their
own hands. This report is about the machinery for deciding it, and about the
much narrower question of what that machinery has actually been shown to do.

Concretely, the input is 30\,fps webcam video of a person solving a
$3\times3\times3$ cube. The output we want is the sequence of quarter turns
they performed, in order --- a string over a 12-symbol alphabet
$\{\wca{U},\wca{U'},\wca{D},\wca{D'},\wca{L},\wca{L'},\wca{R},\wca{R'},
\wca{F},\wca{F'},\wca{B},\wca{B'}\}$ --- together with a verdict on whether
that sequence is consistent with a genuine solve.

\subsection{Why this is not a standard video-recognition problem}

Three properties make it its own thing, and each one drove a design choice
that this report will defend or retract.

\paragraph{The events are dense, brief, and never separated by rest.}
A competent solve is 80--150 quarter turns in 30--60 seconds. In the corpus
used here the median rate is \MedianTPS\ turns per second, and
\CrowdedPctHoldout\% of adjacent turn pairs in the held-out solves land
within two frames of each other at 30\,fps. There is no still moment between
moves to segment on. Any method whose first step is ``find the boundaries of
each move'' is fighting the data, and \cref{sec:peakpick} shows measurably
that it loses.

\paragraph{The label is cheap but subtly wrong.}
A Bluetooth smart cube reports every turn it feels, timestamped. That is an
enormous gift: free, dense supervision for a video model, with no human
annotation at all. But the cube reports turns \emph{relative to its own
core}, and a middle-slice turn rotates that core. After an \wca{M} the
cube's idea of ``up'' has moved and the camera's has not, so every later
label names a face that the camera never saw turn. This is a small effect
(\SliceMoves\ moves corpus-wide) but it is a \emph{systematic} one, it was
invisible to the repository's own consistency check for weeks, and
\cref{sec:labelframe} shows it is still present in every evaluation harness
in the repository other than the ones written for this report.

\paragraph{The output space has algebraic structure, and the endpoints are
known.} This is the property that makes the problem tractable at all. The
$4.3\times 10^{19}$ cube states form a group; a move sequence is a word in
that group; and we know both endpoints of the word --- the scramble (issued
by the server, so known exactly) and the end state (solved, by claim). A
recogniser with 90\% per-move accuracy over a 120-move solve gets the whole
word right essentially never ($0.9^{120} \approx 3\times10^{-6}$). A decoder
that searches for the \emph{cheapest edit of the recognised word that
actually reaches solved} converts an impossible sequence-level problem into
a tractable constrained one. \cref{sec:reconstruct} develops this and
\cref{sec:decode-results} measures how much it buys, which is less than the
framing suggests and for a reason worth understanding.

\subsection{The pipeline in one paragraph}

Frames are cropped to the cube by a YOLOv8 detector whose boxes are
median-smoothed and interpolated across time, resized to $96\times96$, and
stacked as four channels (RGB plus a signed luma difference from the
previous frame). A shared-weight CNN encodes each frame to 128 dimensions; a
four-block dilated temporal convolutional network with a 61-frame receptive
field mixes them; a $1\times1$ convolution emits, per frame, a distribution
over 12 quarter turns plus a CTC blank. The network is trained with a
connectionist temporal classification loss, so it is never told \emph{when} a
move happened, only which moves happened in what order. At inference a CTC
prefix beam search reads a move string off the posteriorgram, and a second
beam search over group elements edits that string --- substituting,
deleting, inserting --- until it is a word that carries the known scramble to
solved, at minimum cost under a probabilistic cost model.
\cref{fig:pipeline} is the whole thing.

\subsection{Contributions of this report}

This is a study document, so the contributions are as much about
\emph{measurement discipline} as about method.

\begin{enumerate}[leftmargin=1.4em,itemsep=3pt]
\item \textbf{A doubled, strictly-defined holdout.} The repository's headline
numbers were computed on eight sessions, two of which were the checkpoints'
own validation sessions. I prepared \NewlyPrepared\ additional never-seen
sessions and defined the holdout programmatically from each checkpoint's own
recorded training and validation lists. The honest held-out set is now
\NHoldoutSolves\ free solves and \NHoldoutScrambles\ prescribed scrambles.
\item \textbf{Corrected ground truth.} All evaluation here scores against
the camera-frame move name, not the cube-frame one the smart cube reports
(\cref{sec:labelframe}).
\item \textbf{A single-holdout ablation ladder} over all five checkpoint
generations (\cref{sec:ablation}), paired per session and tested, which is
the first time the training changes have been comparable to each other. Two
findings fall out. Speed augmentation --- introduced for an anticheat
calibration reason and recorded in the repository as costing nothing --- is
the largest and most significant accuracy intervention on record
(\DSpeedAug\ points, $p$ \DSpeedAugP). And the CTC migration, the project's
headline architectural change, is \DCtcVsPeak\ points with a confidence
interval spanning zero.
\item \textbf{A proximity ladder} (\cref{sec:ladder}) separating
memorisation from generalisation: \LadderTrain\% on trained sessions,
\LadderVal\% on same-day validation, \LadderHeld\% on genuinely unseen days.
\item \textbf{A negative result on the group-theoretic decode}
(\cref{sec:decode-results}): on this holdout it is worth \DecodeGain\ points
and verifies \DecodeVerified\ of \DecodeTrials\ session-seed pairs. The
repository reports ${\sim}+1.2$ points; the difference is entirely the
holdout, and \cref{sec:decode-limits} explains why a two-point-harder test
set collapses the contribution to zero.
\item \textbf{A falsifiability sweep} (\cref{sec:falsifiability}) --- every
claim re-decoded against deliberately false start states --- together with
an argument for why its result here is \emph{uninformative} rather than
reassuring, and a constructive bound showing a cheap wrong story provably
exists even where the search does not find it.
\item \textbf{A catalogue of negative results} (\cref{sec:negatives}). Nine
interventions were measured and found flat or harmful. They are reported
because a reader deciding what to try next needs them more than they need
the positive results.
\end{enumerate}

\subsection{What this report does not claim}

Stated here so it does not have to be re-hedged in every section.

\begin{keybox}
\textbf{One solver, one room, one cube.} Every session in the corpus is the
same person, in one of two rooms, with the same cube. Nothing here
establishes that the model generalises across people, cubes, or
environments; \cref{sec:ladder} shows it degrades measurably across
\emph{days} within one person, which is a lower bound on how much it would
degrade across people.

\smallskip
\textbf{The attack side is a proxy, not an attack.} No adversary has ever
tried to beat this system. The ``caught'' numbers in
\cref{sec:anticheat-results} come from replaying prescribed scrambles, which
are structurally similar to the attack of interest but were not produced by
anybody trying to cheat.

\smallskip
\textbf{Verification and accuracy are different products.} The
group-theoretic decoder makes the transcript more accurate. It does
\emph{not} make a \texttt{VERIFIED} verdict trustworthy, and
\cref{sec:falsifiability} shows exactly where it stops being evidence.
\end{keybox}
